\documentclass[crop=false]{standalone}
% --- ПРЕАМБУЛА ДЛЯ ПОДДЕРЖКИ РУССКОГО ЯЗЫКА И МАТЕМАТИКИ ---
\usepackage[T2A]{fontenc}
\usepackage[utf8]{inputenc}
\usepackage[russian]{babel}
\usepackage{amsmath}
\usepackage{geometry}
\usepackage{amssymb}
\usepackage{graphicx}
\usepackage{float}
\usepackage{import}
\geometry{a4paper, top=2cm, bottom=2cm, left=2cm, right=2cm}
\newcommand{\sidecomment}[1]{\marginpar{\raggedright\scriptsize\itshape #1}}
\newcommand{\iu}{\mathrm{i}}

\newcommand{\Def}{%
    \text{Def}%
    \hspace{-0.85em}% Сдвиг назад ровно на середину (подберите под шрифт)
    \raisebox{0.1ex}{/}% Черта
    \hspace{1em}% Возврат к концу слова + небольшой отступ
}

\renewcommand{\Re}{\operatorname{Re}}
\renewcommand{\Im}{\operatorname{Im}}

\ifstandalone
    \newcommand{\lecturebreak}{} % В отдельном файле лекции разрыв не нужен
\else
    \newcommand{\lecturebreak}{\clearpage} % В полной книге - начинаем с новой страницы
\fi

\newcounter{lecturenumber} % Создаем счетчик для нумерации лекций
\newcommand{\lecture}[2]{
    \lecturebreak
    \refstepcounter{lecturenumber} % Увеличиваем счетчик на 1
    \par\vspace{2em}\noindent % Отступ сверху и убираем абзацный отступ
    {\Large\bfseries % Делаем заголовок крупным и жирным
    Лекция \thelecturenumber: #1 % Выводим "Лекция N: Тема"
    \hfill % Растягиваем пробел, чтобы дата была справа
    \large\textit{#2}} % Выводим дату курсивом
    \par\vspace{1em} % Небольшой отступ снизу
    \addcontentsline{toc}{section}{Лекция \thelecturenumber: #1} % Добавляем в оглавление
    \par
}

\newcommand{\TwoCols}[2]{%
  \par\noindent % Начать с новой строки без отступа
  \begin{minipage}[t]{0.48\textwidth}
    #1
  \end{minipage}% <-- Важный '%' для удаления пробела
  \hfill % Растягивающийся пробел между колонками
  \begin{minipage}[t]{0.48\textwidth}
    #2
  \end{minipage}% <-- Важный '%' для удаления пробела
  \par%
}

\pagestyle{plain} % Включаем нумерацию страниц\
\begin{document}
\lecture{}{20.03}
\Def Cистемы сил называются эквивалентными если замена одной на другую не изменяет движения
\[\{\bar F_k\} \sim \{\bar F_k'\}\]

\begin{theorem}[Критерий экв-ти систем сил]
    Системы сил эквивалентны
    
    $\iff$

    элементарные работы соврешаемые этими силами на виртуальных перемещения равныэ
    

    \[\{\bar F_k\} \sim \{\bar F_k'\} \iff \delta A  = \delta A'\]
\end{theorem}
\begin{proof}
    \[\sum_{k=1}^{N} (\bar F_l - m_k \bar w_k) \delta \bar r_k = 0 \qquad \underbracket{\sum_{k}^{N} \bar F_k \delta \bar r_k}_{\delta A} = \sum_{k}^{N} m_k \bar w_k \delta \bar r_k\]
    \[\delta A' = \sum_k^N m_k \bar w_k' \delta \bar r_k \]
    Системы сил эквиваленты, то есть $\bar w_k = \bar w_k'$, значит $\delta A = \delta A'$

    В обратную сторону доказывается также $\delta A = \delta A'$, значит $\sum_k^N m_k \bar w_k = \sum_k^N m_k \bar w_k'$, к тому же $\bar w_k$ и $\bar w_k'$ связаны одними и теми же соотношения,  тогда $\bar w_k = \bar w_k'$
\end{proof}

\noindent \keypoint Для голономных систем:
\[\delta A = \sum_i^n Q_i \delta q_i \qquad \delta A' \sum_i^n Q_i' \delta q_i\]
\[\delta A = \delta A' \iff \sum_i^n (Q_i - Q_i') \delta q_i = 0\]
В голономных системах $\delta q_i$ независимы, тогда $Q_i = Q_i'$

То есть эквивалентные системы связей приводят к одинковым обобщённым силам

\section{Абсолютно твёрдое тело}
Твёрдое тело это склерономное голономная система, то есть виртуальные связи являются возможными
\[
\begin{aligned} 
    \delta A &= \sum_k^N \bar F_k \delta \bar r_k = \sum_k^N \bar F_k \Delta \bar r_k = \sum_k^N \bar F_k \bar v_k^* \Delta t = \sum_k^N \bar F_k (\bar v_P^* + \bar \omega \times (\bar r_k - \bar r_P)) \Delta t = \sum_k^N \bar F_k \bar v_P^* \Delta t + \sum_k^N \left( \bar F_k ,\bar \omega \times (\bar r_k - \bar r_p)  \right)\Delta t  \\
    &= \bar R \bar v_P^* \Delta t + \sum_k^N (\bar \omega, (\bar r_k - \bar r_P) \times \bar F_k) \Delta t = \bar R \bar v_P^* \Delta t + \left(\omega , \sum_k^N (\bar r_k - \bar r_P) \times \bar F_k \right) \\
    &= \bar R \bar v_P^* \Delta t + \bar \omega^* \bar M_P \Delta t 
\end{aligned}
\] % fixme
\[
\begin{cases}
\delta A = \bar R \bar v_P^* \Delta t + \bar M_P \bar w^* \Delta t \\ \delta A' = \bar R' \bar v_P^* \Delta t + \bar M_P \bar \omega^* \Delta t 
\end{cases} \implies \boxed{(\bar R - \bar R') \bar v_P^* \Delta t + (\bar M_P - \bar M_P') \bar \omega^* \Delta t = 0}
\]
из того, что возможны случаи $\bar v_P^* \neq 0 \& \bar w_P = 0$ и $\bar v_P^* = 0 \& \bar w_P \neq 0$

это критерий эквивалентности систем сил, приложенных к твёрдому телу
\[\boxed{\begin{matrix}\bar R = \bar R' \\ \bar M_P = \bar M_P'\end{matrix}}\]

\keypoint из принципа вертуальных перемещений 
\[\delta A = 0 \qquad \bar R v^* \Delta t + \bar M_P \bar \omega^* \Delta t = 0 \iff \begin{cases}
    \bar R = 0 \\ 
    \bar M_P = 0
\end{cases}\]

Заметим, что силы можно распределять как удобно % fixme фото в телефоне

\section{}
\Def Линия действия сил — прямая коллинеарная силе и проходящая через точку, к которой данная сила прилоежена

Смещение силы, вдоль её линии действия приводит к эквивалентной системе
\begin{proof}
    упр, главный вектор не меняется, главный момент не меняется
\end{proof}

\keypoint
\Def Равнодействующая сила системе сил, приложенной к твёрдому телу — сила, эквивалентная данной система
\[\bar F^* \sim \{\bar F_1 , \dots , \bar F_N\}\]
$\bar r^*$ — точка приложения равнодействующей силы

\begin{theorem}[Вариньона]
Если система сил имеет равнодействующую, то она равна главному вектору и её момент относительно произвольной точки $P$ равен главному моменту сил системы относительно этой же точки
\[\bar F^* = \bar R \qquad \bar m_p(\bar F^*) = \bar M_P \]
\end{theorem}
Следствия:
1. твердое тело не может находиться равновесии под действием одной ненулевой силы

2.  если тело в равновесии под действием двух сил, то их линии действия совпадают, а сами они противополжны

3. Теорема о трёх силах
если твердое тело находится в равновесии под действием трёх сил, причём линии действия $\bar F_1$ и $\bar F_2$ пересекаются, то третья сила лежит в той же плоскости, что и $\bar F_1$ $\bar F_2$ и её линия действяя пересекается с линиями действиями других двух в той же точке
\[(\bar F_1, \bar F_2, \bar F_3 \qquad \Gamma_1 \cap \Gamma_2 = P) \implies \bar F_3 \in (\bar F_1 , \bar F_2) \quad \& \quad \Gamma_3 \cap \Gamma_1 \cap \Gamma_2 = P\]
4. Система сил, линии действяи которых пересекаются в одной точке имеет равнодействующую

Отличия равнодействующей от главного вектора:

равнодействующая не всегда существует, у равнодействующей есть точка приложения


\section{Система параллельных сил}

\Def Система параллельных сил — система сил, линии действия которых параллельных
\[\bar F_k = \bar n \Phi_k\]
$\Phi_k$ — не модуль силы, скорее проекция на $\bar n$

удобно выбирать $\bar n$, совпадающий по направлению с главным вектором
\[\bar R = \sum_k^N \bar F_k = \sum_k^N \bar \Phi_k = \bar n \sum_k^N \Phi_k \]
\end{document}